\section{第41届复赛第7题：超导环磁通量子化与无相互作用电子}
\subsection{题干}
题目以半径为 $R$ 的一维圆环为对象，电子波函数为 $\psi(x)=\exp(ipx/\hbar)$，其中 $x=R\theta$。先由波函数单值性求电子动量的可能取值，再向圆环通入磁通 $\Phi$，讨论单电子能级、多电子最低总能量的周期性、奇偶电子数导致的平均能量差异，以及双环体系中的配对条件。
\begin{figure}[htbp]\centering\begin{tikzpicture}[>=Stealth,scale=1]
  \draw[line width=1.35pt] (0,0) ellipse (2.35 and .82);
  \fill (-1.73,-.51) circle (2.1pt);
  \node[below left] at (-1.73,-.51) {$e^-$};
  \node[below] at (0,-1.22) {原题图7a：一维圆环上的电子};
\end{tikzpicture}
\caption{第41届复赛第7题原题图7a的TikZ重绘。}\end{figure}
\begin{figure}[htbp]\centering\begin{tikzpicture}[>=Stealth,scale=1]
  \draw[line width=1.35pt] (0,0) ellipse (2.35 and .82);
  \draw[->,line width=1.05pt] (0,-1.25) -- (0,1.55) node[above] {$\Phi$};
  \fill (-1.73,-.51) circle (2.1pt);
  \node[below left] at (-1.73,-.51) {$e^-$};
  \node[below] at (0,-1.62) {原题图7b：穿过圆环的磁通};
\end{tikzpicture}
\caption{第41届复赛第7题原题图7b的TikZ重绘。}\end{figure}
\subsection{参考解答整理}
圆环闭合使 $\psi(x+2\pi R)=\psi(x)$，因此动量量子化为
\[p_n=\frac{n\hbar}{R},\qquad n\in\mathbb Z.\]
通入磁通后，$p'_n=p_n+eA$，且 $\Phi=2\pi RA$，故单电子能量为
\[E_n(\Phi)=\frac{\hbar^2}{2mR^2}\left(n+\frac{e\Phi}{h}\right)^2.\]
多电子系统取互不相同的动量量子数并使总能量最小。将磁通平移一个 $h/e$ 时，可用 $n\mapsto n+1$ 保持求和范围不变，因而
\[E_{\mathrm{tot}}\left(\Phi+\frac{h}{e}\right)=E_{\mathrm{tot}}(\Phi).\]
当 $N=2M+1$ 为奇数且 $-h/(2e)\leq\Phi<h/(2e)$ 时，平均能量可写成
\[\bar E(\Phi)=\frac{e^2\Phi^2}{4\pi^2}\frac{1}{2mR^2}+\frac{\hbar^2}{2mR^2}\frac{M(M+1)}{3}.\]
故奇数电子数时 $\Phi=0$ 处的平均能量低于 $|\Phi|=h/(2e)$ 处；偶数电子数时结论相反，平均能量在 $|\Phi|=h/(2e)$ 附近取较小值。

最后引入配对约束：动量为 $p'$ 和 $-p'$ 的电子必须同时存在或同时不存在。答案表明，两个环电子数均为奇数时，电子数差最小且 $\Phi=0$ 的构型能量最低；两个环电子数均为偶数时，对应最低点在 $\lvert\Phi\rvert=h/(2e)$。对于题设的两种双环构型，常数项分别为
\[C_1=\frac{M(M+1)}{3},\qquad C_2=\frac{M(M+1)}{3}+\frac14,\]
所以 $\lim_{M\to\infty}C_1/C_2=1$，两种构型的电子平均能量之比也趋于 $1$。
\subsection{转写依据}
题干依据 \texttt{CPHO41-Q07\_transcription\_v1.md}；答案公式依据答案第16至19页和 \texttt{CPHO41-Q07\_formula\_check\_v1.csv}。图7a、图7b均由原题图重绘，不嵌入原始页面截图。
